% =============================================================================
% Chapitre 2 — SQL analytique avancé
% =============================================================================
\chapter{SQL analytique avancé : CTE et fonctions de fenêtre}

\epigraphe{Les données ne parlent pas d'elles-mêmes : il faut savoir les interroger.}

La base est construite et remplie. Il faut maintenant la faire parler : combien de prêts sont en
retard, combien la banque doit provisionner, si le modèle classe toujours bien les clients, comment
évoluent les encaissements. Ces questions demandent des outils plus puissants qu'un simple
\code{GROUP BY} : les \textbf{CTE}, les \textbf{vues}, les \textbf{agrégats conditionnels} et surtout
les \textbf{fonctions de fenêtre}, les notions SQL les plus testées en entretien data.

\begin{observation}[Tes fichiers pour ce chapitre]
\begin{itemize}
\item Ton code : \code{sql/03_chapitre2.sql}. Le corrigé complet : \code{corriges/chapitre2.sql}.
\item Prérequis : le chapitre 1 exécuté. « Aujourd'hui » est la date de référence de la base, le 30/09/2026.
\end{itemize}
\end{observation}

\begin{situation}{Mise en situation --- Le comité des risques de 9 heures}
Demain matin, le directeur des risques présente l'état du portefeuille au comité de direction. Il te
demande quatre chiffres pour 9 heures :
\begin{enumerate}
\item la répartition des prêts par tranche de retard, et la part de l'encours à risque ;
\item le montant des provisions IFRS~9, stage par stage ;
\item la preuve que le modèle de scoring classe toujours correctement les clients ;
\item l'évolution des encaissements sur les douze derniers mois.
\end{enumerate}
Il est 17 heures. Tu as une base, PostgreSQL et ce chapitre.
\end{situation}

% =============================================================================
\section{Cours et notions importantes}

\subsection{L'ordre d'exécution d'une requête}

On écrit une requête dans l'ordre \code{SELECT ... FROM ... WHERE ...}, mais PostgreSQL l'exécute dans un
autre ordre (figure~\ref{fig:ordre}). Le connaître explique la moitié des erreurs de débutant.

\begin{figure}[ht]
\centering
\begin{tikzpicture}[font=\footnotesize, >=Stealth, node distance=0.25cm,
  etape/.style={draw=grisTexte, fill=white, rounded corners=2pt, minimum height=0.8cm, align=center, inner sep=4pt},
  fen/.style={etape, draw=bleuML, fill=bleuFond, line width=0.9pt}]
  \node[etape] (a) {\textbf{1} FROM\\JOIN};
  \node[etape, right=of a] (b) {\textbf{2}\\WHERE};
  \node[etape, right=of b] (c) {\textbf{3} GROUP\\BY};
  \node[etape, right=of c] (d) {\textbf{4}\\HAVING};
  \node[fen, right=of d] (e) {\textbf{5} SELECT\\{\footnotesize + fenêtres}};
  \node[etape, right=of e] (f) {\textbf{6}\\DISTINCT};
  \node[etape, right=of f] (g) {\textbf{7} ORDER\\BY};
  \node[etape, right=of g] (h) {\textbf{8}\\LIMIT};
  \foreach \x/\y in {a/b,b/c,c/d,d/e,e/f,f/g,g/h} \draw[->, thick, grisTexte] (\x) -- (\y);
\end{tikzpicture}
\caption{L'ordre logique d'exécution. Les fonctions de fenêtre sont calculées à l'étape 5, après le
filtrage et le regroupement.}
\label{fig:ordre}
\end{figure}

Trois conséquences pratiques :
\begin{enumerate}
\item un alias défini dans le \code{SELECT} n'existe pas encore dans le \code{WHERE} ;
\item on ne peut pas filtrer sur le résultat d'une fonction de fenêtre dans le \code{WHERE} (elle n'est
pas encore calculée). PostgreSQL n'a pas de clause \code{QUALIFY} comme d'autres moteurs : on calcule
la fenêtre dans une CTE, puis on filtre dans la requête suivante ;
\item une fonction de fenêtre peut s'appliquer à un agrégat : \code{sum(sum(encours)) OVER ()} est
valide, car le \code{GROUP BY} est déjà fait quand la fenêtre est calculée.
\end{enumerate}

\subsection{Les CTE : découper une requête en étapes nommées}

\begin{analogie}{la recette de cuisine}
Une recette ne dit pas « mélange les œufs battus avec la farine tamisée et le beurre fondu cuit à feu
doux » en une seule phrase. Elle découpe : \emph{1. Fais fondre le beurre. 2. Bats les œufs.
3. Mélange.} Chaque étape a un nom, et la suivante l'utilise. Une CTE fait la même chose pour une
requête.
\end{analogie}

\begin{definition}{CTE (Common Table Expression)}
Une sous-requête nommée, définie avec \code{WITH nom AS (...)} avant la requête principale. Chaque CTE
peut utiliser les précédentes. Elle n'existe que le temps de la requête.
\end{definition}

\begin{lstlisting}[style=sql]
WITH impayes AS (                          -- étape 1 : les échéances échues non soldées
    SELECT pret_id, date_echeance
    FROM echeancier_remboursements
    WHERE date_echeance <= DATE '2026-09-30' AND montant_paye < montant_du
),
retards AS (                               -- étape 2 : le retard de chaque prêt
    SELECT pret_id, DATE '2026-09-30' - min(date_echeance) AS jours_retard
    FROM impayes
    GROUP BY pret_id
)
SELECT count(*) FILTER (WHERE jours_retard > 90) AS prets_en_defaut   -- étape 3 : la réponse
FROM retards;                                                         -- 7 088
\end{lstlisting}

\begin{center}\small
\begin{tabularx}{\linewidth}{@{}P{3cm} L L@{}}
\toprule
\textbf{Outil} & \textbf{Durée de vie} & \textbf{Quand l'utiliser} \\
\midrule
Sous-requête & la requête & un calcul court, utilisé une fois \\
CTE (\code{WITH}) & la requête & découper un calcul en étapes lisibles \\
Vue & permanente (seule la définition est stockée) & une règle réutilisée par beaucoup de requêtes \\
Vue matérialisée & permanente, résultat stocké & un calcul lourd, rafraîchi périodiquement
(\code{REFRESH MATERIALIZED VIEW}) \\
Table temporaire & la session & un résultat intermédiaire réutilisé par plusieurs requêtes \\
\bottomrule
\end{tabularx}
\end{center}

\begin{piege}{Piège --- Une CTE n'est pas un résultat stocké}
Depuis PostgreSQL 12, une CTE lue une seule fois est \textbf{fusionnée} dans la requête principale,
comme une sous-requête : l'optimiseur peut y pousser les filtres. Une CTE lue plusieurs fois est
calculée une fois. On peut forcer le comportement avec \code{WITH x AS MATERIALIZED (...)} ou
\code{NOT MATERIALIZED}. Et une CTE ne survit pas à la requête : pour réutiliser un résultat dans
plusieurs requêtes, il faut une vue ou une table temporaire.
\end{piege}

\subsection{Les vues : écrire une règle une seule fois}

Le calcul du retard sera utilisé par presque toutes les analyses de ce chapitre. Le réécrire à chaque
fois, c'est risquer qu'une requête le calcule différemment des autres. Une \textbf{vue} enregistre la
requête sous un nom ; on l'interroge ensuite comme une table :
\begin{lstlisting}[style=sql]
CREATE OR REPLACE VIEW v_situation_prets AS
SELECT ... ;                                          -- la règle, écrite une fois

SELECT count(*) FROM v_situation_prets WHERE en_defaut;   -- utilisée partout
\end{lstlisting}
La vue ne stocke aucune donnée : elle est recalculée à chaque interrogation, donc toujours à jour.

\subsection{Les agrégats conditionnels : FILTER}

Pour compter ou sommer seulement certaines lignes d'un groupe, la clause \code{FILTER} s'ajoute à
n'importe quel agrégat :
\begin{lstlisting}[style=sql]
SELECT count(*)                                         AS nb_prets,
       count(*)     FILTER (WHERE statut_pret = 'en_defaut') AS nb_defauts,
       sum(montant_demande) FILTER (WHERE note_risque = 'A') AS montant_note_a
FROM demandes_pret;
\end{lstlisting}
C'est l'équivalent, plus lisible, de \code{sum(CASE WHEN condition THEN x END)}. On l'utilisera aussi
pour lire les paramètres en une ligne : \code{max(valeur) FILTER (WHERE cle = 'lgd')}.

\subsection{Les fonctions de fenêtre}

\begin{imagine}
Au marathon de Clermont, chaque coureur reçoit deux classements : son rang au général et son rang dans
sa catégorie d'âge. Pour les calculer, on ne « regroupe » pas les coureurs en une ligne par catégorie :
chaque coureur garde \textbf{sa} ligne, avec son temps, et on lui \textbf{ajoute} une information
calculée en regardant les autres coureurs de sa catégorie.
\end{imagine}

\begin{definition}{Fonction de fenêtre}
Une fonction calculée pour chaque ligne à partir d'un ensemble de lignes liées à elle, la
\textbf{fenêtre}, \textbf{sans regrouper les lignes}. Elle se reconnaît au mot-clé \code{OVER}.
\end{definition}

La différence avec \code{GROUP BY} est fondamentale :

\begin{center}
\begin{minipage}[t]{0.44\linewidth}\centering\small
\code{GROUP BY note_risque}\\[3pt]
\begin{tabular}{@{}cc@{}}
\toprule
note & \code{count(*)} \\
\midrule
A & 3 \\
B & 2 \\
\bottomrule
\end{tabular}\\[3pt]
{\footnotesize 5 prêts $\rightarrow$ 2 lignes}
\end{minipage}
\hfill
\begin{minipage}[t]{0.52\linewidth}\centering\small
\code{count(*) OVER (PARTITION BY note_risque)}\\[3pt]
\begin{tabular}{@{}ccc@{}}
\toprule
prêt & note & nb dans la note \\
\midrule
101 & A & 3 \\
102 & A & 3 \\
103 & A & 3 \\
104 & B & 2 \\
105 & B & 2 \\
\bottomrule
\end{tabular}\\[3pt]
{\footnotesize 5 prêts $\rightarrow$ 5 lignes}
\end{minipage}
\end{center}

\subsubsection{L'anatomie d'une fenêtre}

\begin{lstlisting}[style=sql]
fonction(...) OVER (
    PARTITION BY colonnes     -- découpe en groupes indépendants (facultatif)
    ORDER BY colonnes         -- ordonne les lignes dans chaque groupe (facultatif)
    ROWS BETWEEN ... AND ...  -- le cadre : quelles lignes autour de la ligne courante (facultatif)
)
\end{lstlisting}

\begin{center}\small
\begin{tabularx}{\linewidth}{@{}P{3cm} P{5.3cm} L@{}}
\toprule
\textbf{Famille} & \textbf{Fonctions} & \textbf{Usage dans le projet} \\
\midrule
Rangs & \code{ROW_NUMBER}, \code{RANK}, \code{DENSE_RANK}, \code{NTILE(n)} & déciles de risque, plus grosses expositions \\
Décalage & \code{LAG}, \code{LEAD}, \code{FIRST_VALUE}, \code{LAST_VALUE} & évolution d'un mois sur l'autre, matrice de transition \\
Agrégats fenêtrés & \code{SUM}, \code{AVG}, \code{COUNT}, \code{MIN}, \code{MAX} \code{OVER (...)} & cumuls, moyennes mobiles, parts du total \\
\bottomrule
\end{tabularx}
\end{center}

\subsubsection{Les rangs : ROW\_NUMBER, RANK, DENSE\_RANK}

Classons cinq prêts par PD décroissante (exemple illustratif, avec un ex æquo) :

\begin{center}\small
\begin{tabular}{@{}ccccc@{}}
\toprule
prêt & PD & \code{ROW_NUMBER()} & \code{RANK()} & \code{DENSE_RANK()} \\
\midrule
P1 & 0,30 & 1 & 1 & 1 \\
P2 & 0,25 & 2 & 2 & 2 \\
P3 & 0,25 & 3 & 2 & 2 \\
P4 & 0,10 & 4 & \textbf{4} & \textbf{3} \\
P5 & 0,05 & 5 & 5 & 4 \\
\bottomrule
\end{tabular}
\end{center}

\begin{itemize}
\item \code{ROW_NUMBER} numérote sans jamais répéter ;
\item \code{RANK} donne le même rang aux ex æquo, puis \textbf{saute} des rangs (1, 2, 2, 4) : c'est
le classement sportif ;
\item \code{DENSE_RANK} donne le même rang aux ex æquo \textbf{sans sauter} (1, 2, 2, 3).
\end{itemize}

\begin{piege}{Piège --- ROW\_NUMBER et les ex æquo}
Entre P2 et P3, qui reçoit le numéro 2 ? \code{ROW_NUMBER} en choisit un au hasard, et le choix peut
changer d'une exécution à l'autre. Pour un résultat reproductible, ajoute toujours un critère de
départage : \code{ORDER BY probabilite_defaut DESC, pret_id}. Question d'entretien classique : « les
trois plus gros prêts en défaut de chaque note » se fait avec \code{RANK} (les ex æquo sont tous
gardés) ou \code{ROW_NUMBER} (exactement trois lignes), selon ce que veut le métier.
\end{piege}

\subsubsection{NTILE : découper en groupes de même taille}

\code{NTILE(n)} répartit les lignes ordonnées en $n$ groupes aussi égaux que possible, numérotés de 1
à $n$.

\begin{calcul}{NTILE(4) sur 10 prêts}
$10 = 4 \times 2 + 2$ : chaque groupe reçoit 2 prêts, et les $2$ prêts restants vont aux
\textbf{premiers} groupes. Les tailles sont donc 3, 3, 2, 2.

Sur ta base, les déciles de l'échantillon de test : $6\,482 = 10 \times 648 + 2$, donc les déciles 1 et
2 contiennent 649 prêts et les huit autres 648. C'est exactement ce que tu obtiendras dans l'exercice
2.3.
\end{calcul}

\subsubsection{LAG et LEAD : regarder la ligne d'avant ou d'après}

\code{LAG(x)} renvoie la valeur de \code{x} sur la ligne \textbf{précédente} de la fenêtre,
\code{LEAD(x)} sur la ligne \textbf{suivante}. Exemple réel : le retard de paiement du prêt 61, échéance
par échéance.
\begin{lstlisting}[style=sql]
SELECT num_echeance, date_echeance, date_paiement,
       date_paiement - date_echeance                                     AS retard,
       LAG(date_paiement - date_echeance) OVER (ORDER BY num_echeance)   AS retard_precedent
FROM echeancier_remboursements
WHERE pret_id = 61 AND num_echeance BETWEEN 13 AND 19
ORDER BY num_echeance;
\end{lstlisting}

\begin{center}\small
\begin{tabular}{@{}cccrr@{}}
\toprule
\code{num_echeance} & \code{date_echeance} & \code{date_paiement} & \code{retard} & \code{retard_precedent} \\
\midrule
13 & 15/03/2026 & 13/03/2026 & $-2$ & NULL \\
14 & 15/04/2026 & 14/04/2026 & $-1$ & $-2$ \\
15 & 15/05/2026 & 18/05/2026 & 3 & $-1$ \\
16 & 15/06/2026 & 15/06/2026 & 0 & 3 \\
17 & 15/07/2026 & 03/08/2026 & 19 & 0 \\
18 & 15/08/2026 & NULL & NULL & 19 \\
19 & 15/09/2026 & 15/09/2026 & 0 & NULL \\
\bottomrule
\end{tabular}
\end{center}

Trois remarques : la première ligne n'a pas de précédent (NULL) ; un retard négatif est un paiement en
avance ; et l'échéance 18, impayée, a un retard NULL, que la ligne 19 hérite comme « précédent ».
\code{LAG(x, 1, 0)} remplace l'absence de ligne précédente par 0, mais pas une valeur NULL existante.

\subsubsection{Les agrégats fenêtrés et le cadre}

Avec \code{ORDER BY}, un agrégat fenêtré devient \textbf{cumulatif}. Le \textbf{cadre}
(\code{ROWS BETWEEN}) précise exactement quelles lignes entrent dans le calcul
(figure~\ref{fig:cadre}).
\begin{lstlisting}[style=sql]
sum(encaisse) OVER (ORDER BY mois)                                        -- cumul depuis le début
avg(encaisse) OVER (ORDER BY mois ROWS BETWEEN 2 PRECEDING AND CURRENT ROW) -- moyenne mobile 3 mois
sum(encaisse) OVER ()                                                     -- total de toutes les lignes
\end{lstlisting}

\begin{figure}[ht]
\centering
\begin{tikzpicture}[font=\small,
  ligne/.style={draw=grisTexte, minimum width=2.6cm, minimum height=0.55cm, fill=white},
  cour/.style={ligne, fill=bleuFond, draw=bleuML, line width=0.9pt}]
  \foreach \i/\m in {1/janvier,2/février,3/mars,4/avril,5/mai,6/juin}
    \node[ligne] (l\i) at (0,-0.6*\i) {\m};
  \node[cour] at (0,-3.0) {\textbf{mai (ligne courante)}};
  \draw[decorate, decoration={brace, amplitude=5pt}, thick, vertCadre] (1.45,-1.55) -- (1.45,-3.25)
     node[midway, right=7pt, align=left, font=\footnotesize] {\textbf{ROWS BETWEEN 2 PRECEDING}\\\textbf{AND CURRENT ROW} : mars, avril, mai};
  \draw[decorate, decoration={brace, amplitude=5pt, mirror}, thick, orangeEx] (-1.45,-0.35) -- (-1.45,-3.25)
     node[midway, left=7pt, align=right, font=\footnotesize] {\textbf{ORDER BY mois}\\(cumul) : de janvier à mai};
\end{tikzpicture}
\caption{Deux cadres de fenêtre pour la ligne de mai.}
\label{fig:cadre}
\end{figure}

\begin{piege}{Piège --- Le cadre par défaut est RANGE, pas ROWS}
Avec un \code{ORDER BY} et sans cadre explicite, PostgreSQL utilise
\code{RANGE BETWEEN UNBOUNDED PRECEDING AND CURRENT ROW}. \code{RANGE} inclut les lignes
\textbf{ex æquo} avec la ligne courante : si deux lignes ont la même date, le « cumul » de la première
contient déjà la seconde. Pour un cumul ligne à ligne, écris \code{ROWS} explicitement. Même cause,
autre symptôme : \code{LAST_VALUE(x) OVER (ORDER BY ...)} renvoie la ligne courante, pas la dernière,
sauf à écrire \code{ROWS BETWEEN UNBOUNDED PRECEDING AND UNBOUNDED FOLLOWING}.
\end{piege}

Quand plusieurs fonctions partagent la même fenêtre, la clause \code{WINDOW} évite de la répéter :
\begin{lstlisting}[style=sql]
SELECT LAG(retard, 1) OVER w, LAG(retard, 2) OVER w
FROM ...
WINDOW w AS (PARTITION BY pret_id ORDER BY num_echeance);
\end{lstlisting}

\subsection{Les sous-totaux : ROLLUP}

\code{GROUP BY ROLLUP (stage)} produit une ligne par stage \textbf{plus une ligne de total}, où la
colonne \code{stage} vaut NULL. On l'affiche avec \code{coalesce(stage::text, 'Total')}. Avec
plusieurs colonnes, \code{ROLLUP (a, b)} donne les sous-totaux par \code{a} puis le total général ; la
fonction \code{GROUPING(a)} permet de distinguer un vrai NULL d'une ligne de sous-total.

\subsection{Les indicateurs du risque, traduits en SQL}

\subsubsection{Le retard et l'encours de chaque prêt}

La définition du DPD (partie 2) se traduit directement :
\begin{lstlisting}[style=sql]
coalesce(date_ref - min(e.date_echeance)
           FILTER (WHERE e.date_echeance <= date_ref AND e.montant_paye < e.montant_du), 0) AS jours_retard
\end{lstlisting}
Le \code{min} cherche l'échéance la plus ancienne parmi celles qui sont échues et non soldées ; s'il
n'y en a aucune, il vaut NULL, et \code{coalesce} le remplace par 0. L'encours (capital non remboursé)
tient compte des paiements partiels au prorata :
$\sum \text{part\_capital} \times (1 - \text{montant\_payé} / \text{montant\_dû})$.

\subsubsection{L'AUC et le Gini, sans bibliothèque}

\begin{definition}{AUC}
L'AUC (aire sous la courbe ROC) est la probabilité qu'un prêt en défaut, tiré au hasard, ait une PD
\textbf{plus élevée} qu'un prêt sain tiré au hasard. 0,5 = hasard, 1 = classement parfait. Le
\textbf{Gini} vaut $2 \times AUC - 1$.
\end{definition}

Cette définition donne une formule calculable en SQL, celle de la statistique de Mann--Whitney. On
range tous les prêts par PD croissante (rang 1 = la plus petite PD). Avec $n_1$ défauts et $n_0$ prêts
sains :
\[
\boxed{\;AUC = \frac{\sum \text{rangs des défauts} - \dfrac{n_1(n_1+1)}{2}}{n_1 \times n_0}\;}
\]
\textbf{Pourquoi ?} La somme des rangs des défauts compte, pour chaque défaut, le nombre de prêts
classés en dessous de lui (plus lui-même). En retirant $n_1(n_1+1)/2$, on enlève les défauts classés
sous d'autres défauts : il reste exactement le nombre de paires (défaut, sain) bien ordonnées, que l'on
divise par le nombre total de paires $n_1 n_0$.

\begin{calcul}{l'AUC de cinq prêts}
Deux défauts (D1 : PD 0,9 ; D2 : PD 0,4) et trois prêts sains (S1 : 0,6 ; S2 : 0,3 ; S3 : 0,1).
\begin{enumerate}
\item \textbf{Par les paires.} 6 paires (défaut, sain). D1 dépasse S1, S2, S3 : 3 paires bien
ordonnées. D2 dépasse S2 et S3, pas S1 : 2 paires. $AUC = 5/6 = 0{,}833$.
\item \textbf{Par les rangs.} Par PD croissante : S3 = 1, S2 = 2, D2 = 3, S1 = 4, D1 = 5. Somme des rangs
des défauts : $3 + 5 = 8$. $AUC = (8 - 2 \times 3 / 2) / (2 \times 3) = 5/6 = 0{,}833$.
\item Gini $= 2 \times 0{,}833 - 1 = 0{,}667$.
\end{enumerate}
En cas d'ex æquo, chaque prêt reçoit le \textbf{rang moyen} de son groupe.
\end{calcul}

\subsubsection{La courbe de gains, la matrice de transition, l'analyse par génération}

\begin{itemize}
\item \textbf{Courbe de gains} : on classe les prêts en déciles de PD décroissante et on mesure la part
cumulée des défauts « captés » par les premiers déciles. Si refuser les 10~\% les plus risqués évite
45~\% des défauts, le modèle est très discriminant.
\item \textbf{Matrice de transition} (\emph{roll rates}) : pour chaque tranche de retard en fin de mois,
la répartition des prêts dans les tranches le mois suivant. Elle montre la dynamique du risque :
quelle part des prêts à 31--60 jours se régularise, quelle part s'aggrave.
\item \textbf{Analyse par génération} (\emph{vintage}) : on regroupe les prêts par trimestre d'octroi et
on suit leur taux de défaut cumulé après 6, 12, 18, 24 mois d'ancienneté (\emph{months on book}, MOB).
Une génération qui fait défaut plus vite que les précédentes signale un relâchement de la politique
d'octroi.
\end{itemize}

\begin{piege}{Piège --- Comparer des générations qui n'ont pas le même âge}
Une génération accordée il y a 3 mois n'a pas encore eu le temps de faire défaut : son taux de défaut
« après 12 mois » n'existe pas. Si on le calcule quand même, les jeunes générations paraissent
excellentes. Pour chaque horizon de $h$ mois, on ne garde que les prêts \textbf{observés pendant au
moins $h$ mois}.
\end{piege}

\subsection{Performance : filtrer avant de joindre}

Une requête analytique sur 1,1 million de mensualités peut prendre des secondes ou des minutes selon
la façon dont elle est écrite. La règle la plus rentable : \textbf{réduire les données avant de les
joindre}. Dans la matrice de transition (exercice 2.6), on joint chaque fin de mois à chaque prêt et à
ses échéances : en ne gardant d'abord que les échéances qui \emph{peuvent} être impayées pendant la
période étudiée, la requête passe de 7 à 3 secondes. \code{EXPLAIN ANALYZE} permet de vérifier
l'effet de chaque réécriture.

% =============================================================================
\section{Exercices guidés}

\begin{exercice}{2.1}{La situation de chaque prêt à la date d'arrêté}
\begin{enumerate}[label=\alph*)]
\item Crée la vue \code{v_situation_prets} : pour chaque prêt accordé avant la date de référence, son
retard (DPD), son encours, ses impayés échus, le nombre de mois restants et un booléen
\code{en_defaut} (retard supérieur au seuil). Les paramètres sont lus dans \code{parametres_risque}.
\item Donne la répartition des prêts \textbf{non soldés} par tranche de retard (à jour, 1--30, 31--60,
61--90, plus de 90 jours) : nombre, encours, part de l'encours total.
\item Calcule le PAR30 et le PAR90.
\end{enumerate}
\end{exercice}

\begin{correction}{2.1}
\textbf{a) La vue.} Une première CTE lit les paramètres en une ligne (agrégats \code{FILTER}). La
seconde joint les prêts à leurs échéances et calcule tout par prêt. Le \code{CROSS JOIN} ajoute la
ligne des paramètres à chaque ligne : c'est le moyen simple de rendre une constante disponible partout.

\begin{lstlisting}[style=long]
CREATE OR REPLACE VIEW v_situation_prets AS
WITH arrete AS (
    SELECT max(valeur_date) FILTER (WHERE cle = 'date_reference')  AS date_ref,
           max(valeur)      FILTER (WHERE cle = 'seuil_defaut_jours') AS seuil_defaut
    FROM parametres_risque
),
calcul AS (
    SELECT p.pret_id, p.client_id, p.note_risque, p.motif_pret, p.duree_mois, p.date_decision,
           a.date_ref, a.seuil_defaut,
           -- DPD : jours écoulés depuis la plus ancienne échéance échue et non soldée (0 si tout est à jour)
           coalesce(a.date_ref - min(e.date_echeance)
                      FILTER (WHERE e.date_echeance <= a.date_ref AND e.montant_paye < e.montant_du), 0) AS jours_retard,
           -- encours : capital pas encore remboursé (un paiement partiel couvre capital et intérêts au prorata)
           round(sum(e.part_capital * (1 - e.montant_paye / e.montant_du)), 2)                         AS encours,
           coalesce(sum(e.montant_du - e.montant_paye) FILTER (WHERE e.date_echeance <= a.date_ref), 0) AS impayes,
           count(*) FILTER (WHERE e.date_echeance > a.date_ref)                                        AS mois_restants
    FROM demandes_pret p
    JOIN echeancier_remboursements e USING (pret_id)
    CROSS JOIN arrete a
    WHERE p.statut_dossier = 'accorde' AND p.date_decision <= a.date_ref
    GROUP BY p.pret_id, a.date_ref, a.seuil_defaut    -- p.pret_id est la clé : les autres colonnes de p en dépendent
)
SELECT pret_id, client_id, note_risque, motif_pret, duree_mois, date_decision, date_ref,
       jours_retard, encours, impayes, mois_restants,
       jours_retard > seuil_defaut AS en_defaut       -- « plus de 90 jours » (art. 178 CRR)
FROM calcul;
\end{lstlisting}

\textbf{Remarque sur le GROUP BY.} On sélectionne \code{p.note_risque}, \code{p.duree_mois}\dots sans
les mettre dans le \code{GROUP BY}. C'est permis parce que \code{p.pret_id}, clé primaire de
\code{demandes_pret}, y figure : PostgreSQL sait que les autres colonnes de \code{p} en dépendent.

\textbf{b) Les tranches.} La CTE classe chaque prêt dans une tranche (les numéros devant les libellés
donnent l'ordre de tri). \code{sum(sum(encours)) OVER ()} calcule le total de toutes les tranches sans
perdre le détail : c'est une fonction de fenêtre appliquée au résultat du \code{GROUP BY}.

\begin{lstlisting}[style=long]
WITH tranches AS (
    SELECT encours,
           CASE WHEN jours_retard = 0  THEN '1. À jour'
                WHEN jours_retard <= 30 THEN '2. 1 à 30 jours'
                WHEN jours_retard <= 60 THEN '3. 31 à 60 jours'
                WHEN jours_retard <= 90 THEN '4. 61 à 90 jours'
                ELSE                         '5. Plus de 90 jours' END AS tranche
    FROM v_situation_prets
    WHERE encours > 0                                  -- les prêts soldés ne portent plus de risque
)
SELECT tranche, count(*) AS nb_prets, sum(encours) AS encours,
       -- SUM(SUM(...)) OVER () : le total de toutes les tranches, sans perdre le détail
       round(100 * sum(encours) / sum(sum(encours)) OVER (), 2) AS part_encours_pct
FROM tranches
GROUP BY tranche
ORDER BY tranche;
\end{lstlisting}

\textbf{c) PAR30 et PAR90.}
\begin{lstlisting}[style=long]
SELECT sum(encours)                                                               AS encours_total,
       round(100 * sum(encours) FILTER (WHERE jours_retard > 30) / sum(encours), 2) AS par30_pct,
       round(100 * sum(encours) FILTER (WHERE en_defaut)         / sum(encours), 2) AS par90_pct
FROM v_situation_prets
WHERE encours > 0;
\end{lstlisting}

\begin{resultat}
\begin{center}\small
\begin{tabular}{@{}lrrr@{}}
\toprule
Tranche & Prêts & Encours (€) & Part de l'encours \\
\midrule
1. À jour & 13\,480 & 73\,884\,266,09 & 55,15~\% \\
2. 1 à 30 jours & 714 & 3\,948\,710,84 & 2,95~\% \\
3. 31 à 60 jours & 201 & 1\,081\,975,60 & 0,81~\% \\
4. 61 à 90 jours & 18 & 69\,932,39 & 0,05~\% \\
5. Plus de 90 jours & 7\,088 & 54\,978\,554,44 & 41,04~\% \\
\bottomrule
\end{tabular}

\smallskip
Encours total : 133\,963\,439,36~€ \qquad PAR30 : \textbf{41,90~\%} \qquad PAR90 : \textbf{41,04~\%}
\end{center}
Les 10\,909 prêts soldés (encours nul) sont exclus : $13\,480 + 714 + 201 + 18 + 7\,088 = 21\,501$ prêts
vivants.
\end{resultat}

\begin{observation}[Un PAR90 de 41~\% : réaliste ?]
Non : dans une banque de crédit à la consommation, le PAR90 se mesure plutôt en quelques pour cent.
Deux raisons à cet écart, qu'il faut savoir expliquer. D'abord, le jeu de données contient 21,9~\% de
défauts, bien plus qu'un portefeuille réel. Ensuite, un prêt en défaut cesse de payer et garde tout son
encours, alors que les prêts sains s'amortissent et 10\,909 sont déjà soldés : dans l'encours
\emph{restant}, les défauts pèsent donc de plus en plus lourd. Le calcul est juste ; c'est le
portefeuille qui est atypique.
\end{observation}
\end{correction}

\begin{exercice}{2.2}{Provisions IFRS 9 : stage, EAD, PD, LGD et perte attendue}
Pour chaque prêt non soldé, détermine son stage (seuils lus dans les paramètres), son EAD (l'encours),
la PD du modèle champion, la PD utilisée selon le stage ($PD_m = 1-(1-PD)^{m/n}$) et l'ECL. Présente le
nombre de prêts, l'EAD, la PD moyenne, l'ECL et le taux de couverture (ECL / EAD) par stage, avec une
ligne de total. Vérifie enfin qu'aucun prêt n'est oublié faute de score.
\end{exercice}

\begin{correction}{2.2}
Trois CTE, trois étapes : les paramètres ; le stage de chaque prêt ; la PD utilisée. La requête finale
agrège avec \code{ROLLUP}.

\begin{lstlisting}[style=long]
WITH parametres AS (
    SELECT max(valeur) FILTER (WHERE cle = 'lgd')                AS lgd,
           max(valeur) FILTER (WHERE cle = 'seuil_defaut_jours') AS seuil_defaut,
           max(valeur) FILTER (WHERE cle = 'seuil_stage2_jours') AS seuil_stage2
    FROM parametres_risque
),
stages AS (
    SELECT s.pret_id, s.encours AS ead, s.duree_mois,
           greatest(s.mois_restants, 1) AS mois_restants,       -- un prêt échu mais impayé : au moins 1 mois
           ps.probabilite_defaut        AS pd_vie_entiere,
           CASE WHEN s.jours_retard > prm.seuil_defaut THEN 3   -- défaut avéré
                WHEN s.jours_retard > prm.seuil_stage2 THEN 2   -- dégradation significative du risque
                ELSE 1 END AS stage
    FROM v_situation_prets s
    JOIN predictions_scoring_ml ps ON ps.pret_id = s.pret_id AND ps.est_champion
    CROSS JOIN parametres prm
    WHERE s.encours > 0
),
pertes AS (
    -- La PD du modèle porte sur toute la durée du prêt. Avec un risque mensuel constant,
    -- la probabilité de défaut sur m mois vaut 1 − (1 − PD)^(m / durée).
    SELECT st.*,
           CASE stage
               WHEN 3 THEN 1                                                                  -- défaut : PD = 1
               WHEN 2 THEN 1 - (1 - pd_vie_entiere) ^ (mois_restants::numeric / duree_mois)   -- vie restante
               ELSE        1 - (1 - pd_vie_entiere) ^ (least(mois_restants, 12)::numeric / duree_mois)  -- 12 mois
           END AS pd_utilisee
    FROM stages st
)
SELECT coalesce(stage::text, 'Total')                    AS stage_ifrs9,
       count(*)                                           AS nb_prets,
       sum(ead)                                           AS ead,
       round(100 * avg(pd_utilisee), 2)                   AS pd_moyenne_pct,
       round(sum(pd_utilisee * prm.lgd * ead), 2)         AS ecl,
       round(100 * sum(pd_utilisee * prm.lgd * ead) / sum(ead), 2) AS taux_couverture_pct
FROM pertes
CROSS JOIN parametres prm
GROUP BY ROLLUP (stage)                                  -- ROLLUP ajoute la ligne de total
ORDER BY stage;

-- Contrôle qualité : un prêt accordé sans score du modèle champion disparaîtrait du calcul (jointure interne)
SELECT count(*) AS prets_sans_score
FROM v_situation_prets s
WHERE NOT EXISTS (SELECT 1 FROM predictions_scoring_ml ps WHERE ps.pret_id = s.pret_id AND ps.est_champion);
\end{lstlisting}

\textbf{Deux détails qui comptent.} \code{mois_restants::numeric} : sans conversion, la division de
deux entiers serait entière ($12/36 = 0$). \code{greatest(mois_restants, 1)} : un prêt dont toutes les
échéances sont passées mais impayées a 0 mois restant, ce qui annulerait sa PD.

\begin{resultat}
\begin{center}\small
\begin{tabular}{@{}lrrrrr@{}}
\toprule
Stage & Prêts & EAD (€) & PD moyenne & ECL (€) & Couverture \\
\midrule
1 & 14\,194 & 77\,832\,976,93 & 2,20~\% & 856\,830,45 & 1,10~\% \\
2 & 219 & 1\,151\,907,99 & 3,80~\% & 29\,949,20 & 2,60~\% \\
3 & 7\,088 & 54\,978\,554,44 & 100,00~\% & 24\,740\,349,50 & 45,00~\% \\
\textbf{Total} & 21\,501 & 133\,963\,439,36 & 34,45~\% & \textbf{25\,627\,129,15} & 19,13~\% \\
\bottomrule
\end{tabular}
\end{center}
Prêts sans score : \textbf{0}.
\end{resultat}

\textbf{Lecture.} Le stage 3 représente 96,5~\% des provisions ; son taux de couverture vaut exactement
la LGD (45~\%), puisque sa PD vaut 100~\%. Les 14\,194 prêts du stage 1 (58~\% de l'EAD) ne coûtent que
1,10~\% de leur encours. Le stage 2 compte 219 prêts : les $201 + 18$ prêts des tranches 31--60 et
61--90 jours de l'exercice précédent.
\end{correction}

\begin{exercice}{2.3}{Le modèle classe-t-il bien ? Déciles de risque sur l'échantillon test}
Sur les prêts de l'échantillon \textbf{test}, découpe les prédictions du modèle champion en 10 déciles
de PD décroissante (décile 1 = les plus risqués). Pour chaque décile : nombre de prêts, PD moyenne
prédite, taux de défaut observé, nombre de défauts et part cumulée des défauts captés.
\end{exercice}

\begin{correction}{2.3}
\begin{lstlisting}[style=long]
WITH test AS (
    SELECT ps.probabilite_defaut AS pd, s.en_defaut::int AS defaut
    FROM predictions_scoring_ml ps
    JOIN echantillons_modele em USING (pret_id)
    JOIN v_situation_prets s   USING (pret_id)
    WHERE ps.est_champion AND em.echantillon = 'test'    -- jamais sur l'apprentissage : résultats trop beaux
),
deciles AS (
    SELECT pd, defaut, NTILE(10) OVER (ORDER BY pd DESC) AS decile   -- décile 1 = les 10 % les plus risqués
    FROM test
)
SELECT decile,
       count(*)                     AS nb_prets,
       round(100 * avg(pd), 2)      AS pd_predite_pct,
       round(100 * avg(defaut), 2)  AS taux_defaut_observe_pct,
       sum(defaut)                  AS nb_defauts,
       -- courbe de gains : part cumulée des défauts captés en refusant les déciles 1 à k
       round(100.0 * sum(sum(defaut)) OVER (ORDER BY decile) / sum(sum(defaut)) OVER (), 1) AS defauts_captes_cumules_pct
FROM deciles
GROUP BY decile
ORDER BY decile;
\end{lstlisting}

\textbf{Les astuces.} \code{en_defaut::int} convertit le booléen en 0 ou 1 : sa moyenne est un taux,
sa somme un nombre de défauts. La dernière colonne empile deux niveaux : \code{sum(defaut)} agrège
chaque décile, puis \code{sum(...) OVER (ORDER BY decile)} cumule les déciles.

\begin{resultat}
\begin{center}\small
\begin{tabular}{@{}crrrrr@{}}
\toprule
Décile & Prêts & PD prédite & Défaut observé & Défauts & Défauts captés (cumul) \\
\midrule
1 & 649 & 97,55~\% & 100,00~\% & 649 & 45,8~\% \\
2 & 649 & 67,59~\% & 64,41~\% & 418 & 75,2~\% \\
3 & 648 & 20,06~\% & 15,28~\% & 99 & 82,2~\% \\
4 & 648 & 10,97~\% & 12,65~\% & 82 & 88,0~\% \\
5 & 648 & 7,47~\% & 10,03~\% & 65 & 92,6~\% \\
6 & 648 & 5,23~\% & 5,56~\% & 36 & 95,1~\% \\
7 & 648 & 3,66~\% & 5,86~\% & 38 & 97,8~\% \\
8 & 648 & 2,31~\% & 3,55~\% & 23 & 99,4~\% \\
9 & 648 & 1,10~\% & 1,08~\% & 7 & 99,9~\% \\
10 & 648 & 0,24~\% & 0,15~\% & 1 & 100,0~\% \\
\bottomrule
\end{tabular}
\end{center}
\end{resultat}

\begin{center}
\centering
\begin{tikzpicture}
\begin{axis}[ybar, bar width=8pt, width=14.5cm, height=6cm, ymin=0, ymax=105,
  xlabel={Décile de PD (1 = les plus risqués)}, ylabel={\%}, xtick={1,...,10},
  legend style={at={(0.5,1.03)}, anchor=south, legend columns=3, draw=none, font=\small},
  tick label style={font=\footnotesize}, label style={font=\small}, grid=major, grid style={gray!20},
  axis lines*=left, enlarge x limits=0.07]
\addplot[fill=bleuML!40, draw=bleuML] coordinates {(1,97.55) (2,67.59) (3,20.06) (4,10.97) (5,7.47) (6,5.23) (7,3.66) (8,2.31) (9,1.10) (10,0.24)};
\addplot[fill=rouge!70, draw=rouge] coordinates {(1,100) (2,64.41) (3,15.28) (4,12.65) (5,10.03) (6,5.56) (7,5.86) (8,3.55) (9,1.08) (10,0.15)};
\addplot[sharp plot, mark=*, thick, vertCadre] coordinates {(1,45.8) (2,75.2) (3,82.2) (4,88.0) (5,92.6) (6,95.1) (7,97.8) (8,99.4) (9,99.9) (10,100)};
\legend{PD prédite, Taux de défaut observé, Défauts captés (cumul)}
\end{axis}
\end{tikzpicture}
\captionof{figure}{Déciles du modèle champion sur les 6\,482 prêts de test : PD prédite et défaut observé sont
proches décile par décile, et les deux premiers déciles concentrent 75~\% des défauts.}
\label{fig:deciles}
\end{center}

\textbf{Lecture.} Deux qualités distinctes apparaissent. La \textbf{discrimination} : le taux de défaut
baisse régulièrement du décile 1 au décile 10, et refuser les 10~\% les plus risqués éviterait 45,8~\%
des défauts. La \textbf{calibration} : dans chaque décile, la PD prédite est proche du taux observé
(67,6~\% contre 64,4~\% dans le décile 2), ce qui autorise à utiliser la PD comme une vraie probabilité,
par exemple dans le calcul des provisions.
\end{correction}

\begin{exercice}{2.4}{Encaissements mensuels : cumul, moyenne mobile et évolution}
Calcule les sommes encaissées chaque mois depuis octobre 2025 (d'après la date de paiement), puis pour
chaque mois : le cumul depuis octobre, la moyenne mobile sur trois mois et l'évolution en pourcentage
par rapport au mois précédent.
\end{exercice}

\begin{correction}{2.4}
\begin{lstlisting}[style=long]
WITH mensuel AS (
    SELECT date_trunc('month', date_paiement)::date AS mois, sum(montant_paye) AS encaisse
    FROM echeancier_remboursements
    WHERE date_paiement >= DATE '2025-10-01'
    GROUP BY 1
)
SELECT mois, encaisse,
       sum(encaisse) OVER (ORDER BY mois)                                           AS cumul_depuis_octobre,
       round(avg(encaisse) OVER (ORDER BY mois ROWS BETWEEN 2 PRECEDING AND CURRENT ROW), 2) AS moyenne_mobile_3_mois,
       round(100 * (encaisse / LAG(encaisse) OVER (ORDER BY mois) - 1), 1)          AS evolution_pct
FROM mensuel
ORDER BY mois;
\end{lstlisting}

\textbf{Pourquoi \code{date_paiement IS NOT NULL} n'est pas écrit ?} Le filtre \code{>= DATE '2025-10-01'}
élimine déjà les NULL : une comparaison avec NULL n'est jamais vraie. Le cumul utilise ici le cadre par
défaut (\code{RANGE}) sans risque, car chaque mois n'apparaît qu'une fois.

\begin{center}
\centering
\begin{tikzpicture}
\begin{axis}[width=14.5cm, height=5.6cm, ymin=4.4, ymax=5.7,
  xtick={1,...,12}, xticklabels={oct.,nov.,déc.,janv.,févr.,mars,avr.,mai,juin,juil.,août,sept.},
  ylabel={Millions d'euros}, legend style={at={(0.02,0.05)}, anchor=south west, draw=none, font=\small},
  tick label style={font=\footnotesize}, label style={font=\small}, grid=major, grid style={gray!20},
  axis lines*=left, yticklabel style={/pgf/number format/.cd, fixed, precision=1, use comma}]
\addplot[mark=*, thick, bleuML] coordinates {(1,5.511) (2,5.524) (3,5.503) (4,5.412) (5,5.262) (6,5.312)
  (7,5.134) (8,5.085) (9,5.029) (10,5.097) (11,4.931) (12,4.556)};
\addplot[dashed, thick, orangeEx] coordinates {(1,5.511) (2,5.518) (3,5.513) (4,5.480) (5,5.392) (6,5.329)
  (7,5.236) (8,5.177) (9,5.083) (10,5.070) (11,5.019) (12,4.861)};
\legend{Encaissé dans le mois, Moyenne mobile 3 mois}
\end{axis}
\end{tikzpicture}
\captionof{figure}{Encaissements mensuels d'octobre 2025 à septembre 2026 (62,36 millions d'euros au total).}
\label{fig:encaissements}
\end{center}

\textbf{Lecture.} Les encaissements baissent lentement : le portefeuille s'amortit, les prêts en
défaut ne paient plus, et les nouveaux prêts se raréfient (les dernières demandes du jeu simulé datent de fin
juin 2026). La chute de septembre ($-7{,}6~\%$) est en partie un \textbf{effet de bord} : les
mensualités de fin septembre payées avec quelques jours de retard ne seront encaissées qu'en octobre,
après la date d'arrêté. Réflexe à avoir en entretien : \emph{la dernière période d'une série est souvent
incomplète}.
\end{correction}

% =============================================================================
\section{Exercices à faire seul}

\begin{exercice}{2.5}{Champion contre challengers : AUC, Gini, rappel et taux de refus}
Sur l'échantillon test, calcule pour \textbf{chacun des trois modèles} : l'AUC (formule des rangs, avec
rangs moyens en cas d'ex æquo), le Gini, le rappel (part des défauts que le modèle aurait refusés), le
taux de refus et le taux de défaut parmi les prêts qu'il aurait accordés. Le champion mérite-t-il sa
place ?
\end{exercice}

\begin{indice}
\begin{itemize}
\item Le rang par modèle : \code{RANK() OVER (PARTITION BY version_modele ORDER BY pd)}.
\item Le rang moyen d'un groupe d'ex æquo de taille $t$ dont le premier rang est $R$ vaut
$R + (t-1)/2$ ; $t$ s'obtient avec \code{count(*) OVER (PARTITION BY version_modele, pd)}.
\item Puis un \code{GROUP BY version_modele} avec des agrégats \code{FILTER}. Attention aux divisions
entières.
\end{itemize}
\end{indice}

\begin{exercice}{2.6}{La matrice de transition entre fins de mois, avec LEAD}
Pour chaque fin de mois de septembre 2025 à septembre 2026, calcule le retard de chaque prêt
\textbf{à cette date} (une échéance est impayée à la date $D$ si elle est échue à $D$ et n'était pas
soldée à $D$). Classe les prêts en tranches, puis donne, pour chaque tranche de départ, la répartition
(en~\%) des prêts dans les tranches \textbf{le mois suivant}. Écarte les prêts terminés et soldés.
\end{exercice}

\begin{indice}
\begin{itemize}
\item Les fins de mois : \code{generate_series} sur les premiers jours de mois, puis
\code{+ INTERVAL '1 month - 1 day'}.
\item « Non soldée à $D$ » : \code{date_paiement IS NULL OR date_paiement > D}.
\item Un prêt sans impayé n'a aucune échéance à joindre : il faut un \code{LEFT JOIN} pour le garder
(avec un retard de 0).
\item La tranche du mois suivant : \code{LEAD(tranche) OVER (PARTITION BY pret_id ORDER BY date_arrete)}.
\item Performance : filtre les échéances inutiles \emph{avant} la jointure.
\end{itemize}
\end{indice}

\begin{exercice}{2.7}{L'analyse par génération (vintage)}
Regroupe les prêts accordés par trimestre d'octroi. Pour chaque génération, calcule le taux de défaut
cumulé après 6, 12, 18 et 24 mois (MOB). La date d'entrée en défaut d'un prêt est la date d'échéance de
sa première mensualité restée impayée plus de 90 jours, plus 91 jours. Ajoute l'écart du MOB 12 avec
la génération précédente. Que remarques-tu sur les générations récentes ?
\end{exercice}

\begin{indice}
\begin{itemize}
\item Une échéance a été impayée plus de 90 jours si
\code{coalesce(date_paiement, date_ref) - date_echeance > 90}.
\item \code{date_trunc('quarter', date_decision)} donne le trimestre.
\item Les horizons : \code{CROSS JOIN (VALUES (6), (12), (18), (24)) AS h (mois)}, en ne gardant que
les prêts observés au moins \code{h.mois} mois (relis le piège des générations).
\end{itemize}
\end{indice}

\begin{exercice}{2.8}{Classements et signal précoce}
\begin{enumerate}[label=\alph*)]
\item Les trois plus grosses expositions en défaut de chaque note de risque, avec le nom du client et
la part de chacune dans l'encours en défaut de sa note.
\item Le \textbf{signal précoce} : les 20 prêts encore à jour ou presque (retard actuel $\leq 30$ jours,
encours non nul) dont le retard de paiement \textbf{augmente} sur chacune de leurs trois dernières
échéances payées, la dernière étant payée en retard. Trie par retard décroissant puis par encours.
\end{enumerate}
\end{exercice}

\begin{indice}
\begin{itemize}
\item a) \code{RANK() OVER (PARTITION BY note_risque ORDER BY encours DESC)} dans une CTE, puis
\code{WHERE rang <= 3} dans la requête suivante.
\item b) \code{LAG(retard, 1)} et \code{LAG(retard, 2)} sur une même fenêtre (clause \code{WINDOW}), et
\code{ROW_NUMBER() OVER (PARTITION BY pret_id ORDER BY num_echeance DESC) = 1} pour la dernière échéance
payée.
\end{itemize}
\end{indice}

% =============================================================================
\section{Synthèse pour les entretiens}

\begin{resume}
\begin{enumerate}
\item \textbf{Ordre d'exécution} : FROM, WHERE, GROUP BY, HAVING, SELECT (fenêtres), DISTINCT, ORDER BY,
LIMIT. On filtre une fenêtre dans une requête englobante (pas de \code{QUALIFY} dans PostgreSQL).
\item \textbf{CTE} pour découper, \textbf{vue} pour réutiliser une règle, \textbf{FILTER} pour les
agrégats conditionnels.
\item \textbf{Fenêtres} : elles gardent les lignes. \code{PARTITION BY} découpe, \code{ORDER BY} ordonne,
le cadre (\code{ROWS}) délimite.
\item \textbf{Rangs} : \code{ROW_NUMBER} (unique), \code{RANK} (saute), \code{DENSE_RANK} (ne saute
pas), \code{NTILE} (groupes égaux, les premiers reçoivent le reste).
\item \textbf{LAG / LEAD} : ligne précédente / suivante ; NULL au bord.
\item \textbf{Cadre par défaut} : \code{RANGE} avec les ex æquo ; écrire \code{ROWS} pour un cumul ligne à
ligne.
\item \textbf{Indicateurs} : DPD, PAR, ECL par stage avec \code{ROLLUP}, déciles et courbe de gains, AUC
par les rangs, matrice de transition, générations observées assez longtemps.
\end{enumerate}
\end{resume}

\begin{entretien}
\begin{itemize}
\item « Le retard est calculé dans une vue unique, \code{v_situation_prets} : toutes les analyses
utilisent la même règle, celle de l'article 178 du CRR. »
\item « Les provisions IFRS~9 tiennent en une requête de trois CTE : 25,6 millions d'euros, dont 96,5~\%
en stage 3. Le \code{ROLLUP} donne le total, et un contrôle vérifie qu'aucun prêt n'est sans score. »
\item « J'ai calculé l'AUC directement en SQL par la statistique de Mann--Whitney, avec des rangs
moyens pour les ex æquo : 0,923 pour le XGBoost contre 0,876 pour la logistique, sur l'échantillon
test uniquement. »
\item « Les déciles montrent à la fois la discrimination, 75~\% des défauts dans les deux premiers
déciles, et la calibration : la PD prédite suit le taux observé. »
\end{itemize}
\end{entretien}

\subsection*{Questions fréquentes en entretien}

\begin{description}[leftmargin=0pt, itemsep=0.5em]
\item[Différence entre WHERE et HAVING ?] \code{WHERE} filtre les lignes avant le regroupement,
\code{HAVING} filtre les groupes après.
\item[RANK, DENSE\_RANK, ROW\_NUMBER ?] Voir le tableau des rangs : le traitement des ex æquo les
distingue.
\item[Comment obtenir le 2\textsuperscript{e} plus gros prêt de chaque client ?]
\code{DENSE_RANK() OVER (PARTITION BY client_id ORDER BY montant DESC)} dans une CTE, puis
\code{WHERE rang = 2}.
\item[Comment calculer une évolution d'un mois sur l'autre ?] \code{LAG(valeur) OVER (ORDER BY mois)},
puis \code{valeur / LAG(...) - 1}.
\item[Une CTE est-elle plus rapide qu'une sous-requête ?] Pas en soi : depuis PostgreSQL 12, elles sont
généralement optimisées de la même façon. La CTE gagne en lisibilité.
\end{description}
